\documentclass{beamer}

%\usepackage[table]{xcolor}
\mode<presentation> {
  \usetheme{Boadilla}
%  \usetheme{Pittsburgh}
%\usefonttheme[2]{sans}
\renewcommand{\familydefault}{cmss}
%\usepackage{lmodern}
%\usepackage[T1]{fontenc}
%\usepackage{palatino}
%\usepackage{cmbright}
  \setbeamercovered{transparent}
\useinnertheme{rectangles}
}
%\usepackage{normalem}{ulem}
%\usepackage{colortbl, textcomp}
\setbeamercolor{normal text}{fg = black}
\setbeamercolor{structure}{fg = blue}
\definecolor{trial}{cmyk}{1,0,0, 0}
\definecolor{trial2}{cmyk}{0.00,0,1, 0}
\definecolor{darkgreen}{rgb}{0,.4, 0.1}
\definecolor{mygray}{gray}{0.8}
\definecolor{myblack}{rgb}{0, 0, 0}
\usepackage{array}
\beamertemplatesolidbackgroundcolor{white}  \setbeamercolor{alerted
text}{fg=red}

\setbeamertemplate{caption}[numbered]\newcounter{mylastframe}

\usepackage{color}
\usepackage{tikz}
\usetikzlibrary{arrows}
\usepackage{colortbl}
\usepackage{graphicx}
\usepackage{epsfig}

%\usepackage[usenames, dvipsnames]{color}
%\setbeamertemplate{caption}[numbered]\newcounter{mylastframe}c
%\newcolumntype{Y}{\columncolor[cmyk]{0, 0, 1, 0}\raggedright}
%\newcolumntype{C}{\columncolor[cmyk]{1, 0, 0, 0}\raggedright}
%\newcolumntype{G}{\columncolor[rgb]{0, 1, 0}\raggedright}
%\newcolumntype{R}{\columncolor[rgb]{1, 0, 0}\raggedright}

%\begin{beamerboxesrounded}[upper=uppercol,lower=lowercol,shadow=true]{Block}
%$A = B$.
%\end{beamerboxesrounded}}
\renewcommand{\familydefault}{cmss}
%\usepackage[all]{xy}

\usepackage{tikz}
\usepackage{lipsum}

 \newenvironment{changemargin}[3]{%
 \begin{list}{}{%
 \setlength{\topsep}{0pt}%
 \setlength{\leftmargin}{#1}%
 \setlength{\rightmargin}{#2}%
 \setlength{\topmargin}{#3}%
 \setlength{\listparindent}{\parindent}%
 \setlength{\itemindent}{\parindent}%
 \setlength{\parsep}{\parskip}%
 }%
\item[]}{\end{list}}
\usetikzlibrary{arrows}
%\usepackage{palatino}
%\usepackage{eulervm}
\usecolortheme{lily}

\newtheorem{com}{Comment}
\newtheorem{lem} {Lemma}
\newtheorem{prop}{Proposition}
\newtheorem{thm}{Theorem}
\newtheorem{defn}{Definition}
\newtheorem{cor}{Corollary}
\newtheorem{obs}{Observation}
 \numberwithin{equation}{section}


\title[PLSC 43401] % (optional, nur bei langen Titeln nötig)
{Mathematical Foundations of Political Methodology}

\author{Robert Gulotty}
\institute[University of Chicago]{Associate Professor\\Department of Political Science \\  University of Chicago}
\vspace{0.3in}
\date{August 28th, 2017}

\setbeamertemplate{navigation symbols}{}

\begin{document}

\begin{frame}
\titlepage
\end{frame}



\begin{frame}

\begin{itemize}
\item \textcolor{myblack}{Quadratic Form}
\item \textcolor{myblack}{Positive Definiteness}
\item \textcolor{myblack}{Multivariate Optimization}
\end{itemize}

\end{frame}



\begin{frame}

\begin{itemize}
\item \textcolor{myblack}{Quadratic Form}
\item \textcolor{mygray}{Positive Definiteness}
\item \textcolor{mygray}{Multivariate Optimization}
\end{itemize}

\end{frame}



\begin{frame}
\frametitle{Quadratic Form}

\ \\Before we proceed our discussion to multivariate optimization, we introduce some concepts related to it. The first important concept is positive/negative definiteness.

\end{frame}



\begin{frame}
\frametitle{Quadratic Form}

\begin{itemize}
\item A \textit{quadratic form} is a homogenous quadratic polynomial in $n$ variables \pause
\begin{itemize}
\item A \textit{homogeneous polynomial} is a polynomial whose nonzero terms all have the same degree: in words, the sum of the exponents of each term is constant \pause
\item A \textit{homogenous quadratic polynomial} is a polynomial whose terms all have degree 2, see following examples: \pause

\begin{itemize}
\item $a x^2$
\item $ax^2+bxy+cy^2$
\item $a x^2+b y^2+c z^2+d xy+e xz+f yz$ (with coefficients $a, b, c, d, e, f$) \pause
\item $a x_{2} + b y$ is not a homogenous quadratic polynomial because one of the terms doesn't have degree 2
\end{itemize}

\end{itemize}

\end{itemize}

\end{frame}



\begin{frame}
\frametitle{Quadratic Form}

Consider the two variable quadratic form
$$q(x, y)=ax^2+bxy+cy^2$$ \pause

We can re-write it as: $$q(x, y)=\mathbf{z}^T\mathbf{A}\mathbf{z}$$ where $$\mathbf{z}=\begin{bmatrix}
x\\y
\end{bmatrix} \hbox{ and }\mathbf{A}=
\begin{bmatrix}
a & s \\
t& c
\end{bmatrix} 
$$
with $s+t=b$ \pause

\ \\If we impose the condition that $\mathbf{A}$ is symmetric (so that $\mathbf{A}=\mathbf{A}^T$, or $a_{ij}=a_{ji}$ for all $i, j$) then we have $s=t={b\over 2}$ \pause
 
 \ \\There is therefore a unique symmetric matrix $\mathbf{A}$ satisfying $q(x, y)=\mathbf{z}^T\mathbf{A}\mathbf{z}$

\end{frame}



\begin{frame}
\frametitle{Quadratic Form}

This generalizes: \pause

\ \\Let $q(\mathbf{x})$ be a quadratic form in $n$ variables, $x_1, ..., x_n$ \pause

\ \\Then $q(\mathbf{x})$ is equivalent to the expression

$$\mathbf{x}^T\mathbf{A}\mathbf{x}$$ for a unique symmetric $n\times n$ matrix $\mathbf{A}$ and $\mathbf{x}^T=[x_1, x_2, ..., x_n]$ \pause

\ \\Takeaway point: A symmetric $\mathbf{A}$ is uniquely determined by the coefficients in  $q(\mathbf{x})$

\end{frame}



\begin{frame}

\begin{itemize}
\item \textcolor{mygray}{Quadratic Form}
\item \textcolor{myblack}{Positive Definiteness}
\item \textcolor{mygray}{Multivariate Optimization}
\end{itemize}

\end{frame}



\begin{frame}
\frametitle{Positive Definiteness}

\ \\For a given quadratic form $q(\mathbf{x})$ we can ask whether $q(\mathbf{x})>0$ (or less stringently, $q(\mathbf{x})\geq 0$) for every non-zero $\mathbf{x}$ \pause

\ \\If  $q(\mathbf{x})>0$ for every non-zero $\mathbf{x}$ then $q(\mathbf{x})$ is \textit{positive definite} \pause

\begin{itemize}
\item Example $q(\mathbf{x}) = x_1^2+x_2^2+...+x_n^2$
\end{itemize} \pause

\ \\If  $q(\mathbf{x})\geq 0$ for every $\mathbf{x}$ then $q(\mathbf{x})$ is \textit{positive semidefinite} \pause

\begin{itemize}
\item Example $q(\mathbf{x}) = x_1^2+x_2^2+x_3^2-2x_2x_3$ \pause
\item Notice that we can rewrite it as: $q(\mathbf{x}) = x_1^2+(x_2-x_3)^2 \geq 0$ but $q(0, 1, 1)=0$
\end{itemize}

\end{frame}



\begin{frame}
\frametitle{Positive Definiteness}

\ \\If  $q(\mathbf{x})<0$ for every non-zero $\mathbf{x}$ then $q(\mathbf{x})$ is \textit{negative definite} \pause

\ \\If  $q(\mathbf{x})\leq 0$ for every $\mathbf{x}$ then $q(\mathbf{x})$ is \textit{negative semidefinite} \pause

\ \\Most  quadratic forms are neither positive nor negative semidefinite \pause

\ \\Example: $q(x, y)=xy$
\begin{itemize}
\item $q(1, 1)>0$
\item $q(1, -1)<0$
\end{itemize}

\end{frame}



\begin{frame}
\frametitle{Positive Definiteness}

\ \\We can test if a quadratic form is positive (negative) (semi)definite by testing the associated symmetric matrix $\mathbf{A}$ that satisfies $q(\mathbf{x})=\mathbf{x}^T\mathbf{A}\mathbf{x}$ \pause

\ \\If $\mathbf{A}$ is $2\times 2$, and recall, $\mathbf{A}=\begin{bmatrix}a& c\\c& d\end{bmatrix}$, then: \pause

\begin{itemize}
\item $\mathbf{A}$ is PD iff $a, d>0$ and $\det \mathbf{A}>0$ \pause
\item $\mathbf{A}$ is PSD iff  $a, d\geq0$ and $\det \mathbf{A}\geq 0$ \pause
\item $\mathbf{A}$ is ND iff $a, d<0$ and $\det \mathbf{A}>0$ \pause
\item $\mathbf{A}$ is NSD iff  $a, d\leq $ and $\det \mathbf{A}\geq 0$
\end{itemize}

\end{frame}



\begin{frame}

\begin{itemize}
\item \textcolor{mygray}{Quadratic Form}
\item \textcolor{mygray}{Positive Definiteness}
\item \textcolor{myblack}{Multivariate Optimization}
\end{itemize}

\end{frame}



\begin{frame}
\frametitle{Multivariate Optimization}

Now, we can proceed our discussion to optimize multivariate functions \pause
\begin{itemize}
\item[-] Parameters $\boldsymbol{\beta} = (\beta_{1}, \beta_{2}, \hdots, \beta_{n} ) $ such that $f(\boldsymbol{\beta}| \boldsymbol{X}, \boldsymbol{Y})$ is maximized
\item[-] Policy $\boldsymbol{x} \in \Re^{n}$ that maximizes $U(\boldsymbol{x})$
\item[-] Weights $\boldsymbol{\pi} = (\pi_{1}, \pi_{2}, \hdots, \pi_{K})$ such that a weighted average of forecasts $\boldsymbol{f}  =  (f_{1} , f_{2}, \hdots, f_{k})$ have minimum loss 
\begin{eqnarray}
\min_{\boldsymbol{\pi}} & = & - (\sum_{j=1}^{K} \pi_{j} f_{j}  - y ) ^ 2 \nonumber 
\end{eqnarray} 
\end{itemize} \pause

We'll describe analytic and computational approaches to optimization

\end{frame}




\begin{frame}
\frametitle{Multivariate Optimization}

\begin{defn} 
Let $\boldsymbol{x} \in \Re^{n}$ and let $\delta >0$.  Define a \alert{neighborhood} of $\boldsymbol{x}$, $B(\boldsymbol{x}, \delta)$, as the set of points such that, 
\begin{eqnarray}
B(\boldsymbol{x}, \delta) & = & \{ \boldsymbol{y} \in \Re^{n} : ||\boldsymbol{x} - \boldsymbol{y}||< \delta \}  \nonumber 
\end{eqnarray}
\end{defn}

\end{frame}



\begin{frame}
\frametitle{Multivariate Optimization}


\begin{defn} Suppose $f:X \rightarrow R$ with $X \subset R^{n}$.  A vector $\boldsymbol{x}^{*} \in X$ is a \alert{global maximum} if , for all other $\boldsymbol{x} \in X$ 
\begin{eqnarray}
f(\boldsymbol{x}^{*}) & > & f(\boldsymbol{x} ) \nonumber 
\end{eqnarray}

A vector $\boldsymbol{x}^{\text{local}}$ is a \alert{local} maximum if there is a neighborhood around $\boldsymbol{x}^{\text{local}}$, $Q \subset X$ such that, for all $x \in Q$, 
\begin{eqnarray}
f(\boldsymbol{x}^{\text{local} }) & > & f(\boldsymbol{x} ) \nonumber 
\end{eqnarray}
\end{defn}

\end{frame}




%\begin{frame}
%\frametitle{Multivariate Optimization}
%\begin{defn} A set $X\subset R^{n}$ is \alert{compact} if it is closed and bounded
%\end{defn} 

%\begin{thm} \alert{Multivariate Extreme Value Theorem} Suppose $f:X \rightarrow \Re$ be continuous and $X \subset \Re^{n}$ and $X$ compact.  Then $f$ takes on its \alert{maximum} and \alert{minimum} values on $X$.  
%\end{thm}

%We're going to come up with the multivariate equivalent of the \alert{first order} and \alert{second order} conditions now

%\end{frame}




\begin{frame}
\frametitle{Gradient} 

\begin{defn} 
Suppose $f:X \rightarrow \Re^{n}$ with $X \subset \Re^{1}$ is a differentiable function.  Define the gradient vector of $f$ at $\boldsymbol{x}_{0}$, $\nabla f(\boldsymbol{x}_{0})$ as, 
\begin{eqnarray}
\nabla f (\boldsymbol{x}_{0})  & = & \left(\frac{\partial f (\boldsymbol{x}_{0}) }{\partial x_{1} }, \frac{\partial f (\boldsymbol{x}_{0}) }{\partial x_{2} }, \frac{\partial f (\boldsymbol{x}_{0}) }{\partial x_{3} }, \hdots, \frac{\partial f (\boldsymbol{x}_{0}) }{\partial x_{n} } \right) \nonumber 
\end{eqnarray}

\end{defn}



\end{frame}



\begin{frame}
\frametitle{Gradient First Order Condition}

\begin{thm}
Suppose $f:X \rightarrow R^{1}$, $X \subset R^{n}$.  Suppose $\boldsymbol{a} \in X$ is a local extremum.  Then, 
\begin{eqnarray}
\nabla f(\boldsymbol{a}) & = & \boldsymbol{0} \nonumber \\
									& = & (0, 0, \hdots, 0) \nonumber 				
\end{eqnarray}

\end{thm} \pause

\begin{itemize}
\item[-] Proof (intuition): same as one dimensional case (left-hand, right hand), just do it dimension by dimension \pause
\item[-] Critical Values: 
\begin{itemize}
\item[1)] Maximum
\item[2)] Minimum 
\item[3)] Saddle point
\end{itemize}
\item[-] Second Derivative Test!
\end{itemize}

\end{frame}



\begin{frame}
\frametitle{Second Order Conditions: Hessian}


\begin{defn}
Suppose $f:X \rightarrow R^{1}$ , $X \subset R^{n}$, with $f$ a twice differentiable function.  We will define the Hessian matrix as the matrix of second derivatives at $\boldsymbol{x}^{*} \in X$,
\begin{eqnarray}
\boldsymbol{H}(f)(\boldsymbol{x}^{*} )  & = & \begin{pmatrix} 
		\frac{\partial^{2} f }{\partial x_{1} \partial x_{1} } (\boldsymbol{x}^{*} ) & \frac{\partial^{2} f }{\partial x_{1} \partial x_{2} } (\boldsymbol{x}^{*} ) & \hdots & \frac{\partial^{2} f }{\partial x_{1} \partial x_{n} } (\boldsymbol{x}^{*} ) \\
		\frac{\partial^{2} f }{\partial x_{2} \partial x_{1} } (\boldsymbol{x}^{*} ) & \frac{\partial^{2} f }{\partial x_{2} \partial x_{2} } (\boldsymbol{x}^{*} ) & \hdots & \frac{\partial^{2} f }{\partial x_{2} \partial x_{n} } (\boldsymbol{x}^{*} ) \\
		\vdots & \vdots & \ddots & \vdots \\
		\frac{\partial^{2} f }{\partial x_{n} \partial x_{1} } (\boldsymbol{x}^{*} ) & \frac{\partial^{2} f }{\partial x_{n} \partial x_{2} } (\boldsymbol{x}^{*} ) & \hdots & \frac{\partial^{2} f }{\partial x_{n} \partial x_{n} } (\boldsymbol{x}^{*} ) \\
\end{pmatrix} \nonumber 
\end{eqnarray}

\end{defn} \pause	
								
General test $\leadsto$ Two Dimensional Test $\leadsto$ Example

\end{frame}




\begin{frame}
\frametitle{Second Derivative Test} 

Many ways to assess definiteness $\leadsto$ use determinant \pause

\begin{thm} 
Two Dimensional, Second Derivative Test.  Suppose $f:X \rightarrow R$ with $X \subset R^{2}$ and $f$ twice differentiable.  Write the Hessian of $f$ at a critical value $\boldsymbol{a}$, 
\begin{eqnarray}
\boldsymbol{H}(f)(\boldsymbol{a}) & = & \begin{pmatrix} 
	A & B \\
	B & C \\
	\end{pmatrix} \nonumber 
	\end{eqnarray}
Then, we can conduct the second derivative test as:
\begin{itemize}
\item[-] $AC - B^2> 0$ and $A>0$ $\leadsto$ positive definite $\leadsto$ $\boldsymbol{a}$ is a local minimum 
\item[-] $AC - B^2> 0 $ and $A<0$ $\leadsto$ negative definite $\leadsto$ $\boldsymbol{a}$ is a local maximum
\item[-] $AC - B^2<0 $ $\leadsto$ indefinite $\leadsto$ saddle point 
\item[-] $AC- B^2 = 0$ inconclusive
\end{itemize}
	
\end{thm}

\end{frame}



\begin{frame}
\frametitle{Multivariate Recipe}

\begin{itemize}
\item[1)] Calculate gradient \pause
\item[2)] Set equal to zero, solve system of equations \pause
\item[3)] Calculate Hessian \pause
\item[4)] Assess Hessian at critical values \pause
\item[5)] Boundary values?  (if relevant)
\end{itemize}

\end{frame}



\begin{frame}
\frametitle{Example 1: A Simple Optimization Problem}

Suppose $f:R^{2} \rightarrow \Re$ with 
\begin{eqnarray}
f(x_{1}, x_{2}) & = & 3(x_1 + 2)^2  + 4(x_{2}  + 4)^2 \nonumber 
\end{eqnarray} \pause

Calculate gradient \pause

\begin{eqnarray}
\nabla f(\boldsymbol{x}) & = & (6 x_{1} + 12 , 8x_{2} + 32 ) \nonumber \\
\boldsymbol{0} & = & (6 x_{1}^{*} + 12 , 8x_{2}^{*} + 32 ) \nonumber 
\end{eqnarray} \pause

We now solve the system of equations to yield
$x_{1}^{*}  = - 2$ and $x_{2}^{*}  = -4 $ 

\end{frame}



\begin{frame}
\frametitle{Example 1: A Simple Optimization Problem}

\begin{eqnarray}
\textbf{H}(f)(\boldsymbol{x}^{*}) & = & \begin{pmatrix}
6 & 0 \\
0 & 8 \\
\end{pmatrix}\nonumber 
\end{eqnarray} \pause

det$(\textbf{H}(f)(\boldsymbol{x}^{*}))$ = 48 and $6>0$ so $\textbf{H}(f)(\boldsymbol{x}^{*})$ is positive definite. \alert{local minimum}

\end{frame}





\begin{frame}
\frametitle{Example 2: Two Dimensional Ideal Points}
Suppose legislators are considering legislation $\boldsymbol{x} \in R^{2}$.  And suppose legislator $i$ has utility function $U_{i}: R^{2} \rightarrow \Re$, \pause
\begin{eqnarray}
U(\boldsymbol{x})_{i} & = & - (x_{1} - \mu_{1})^2 - (x_{2} - \mu_{2})^2 \nonumber 
\end{eqnarray} \pause
What is legislator $i$'s optimal policy? \pause

$\nabla f(\boldsymbol{x}) = ( -2 (x_{1} - \mu_{1} ) , -2 (x_{2} - \mu_{2} )  )$\\
$\nabla f(\boldsymbol{x}) = \boldsymbol{0} $ \pause
\begin{eqnarray}
- 2(x_{1}^{*} - \mu_{1} ) & = & 0 \nonumber \\
- 2(x_{2}^{*} - \mu_{2} ) & = & 0 \nonumber 
\end{eqnarray}  \pause
Solving yields $x_{1}^{*} = \mu_{1}$ and $x_{2}^{*} = \mu_{2}$.  

\end{frame}



\begin{frame}
\frametitle{Example 2: Two Dimensional Ideal Points} 

\begin{eqnarray}
U(\boldsymbol{x})_{i} & = & - (x_{1} - \mu_{1})^2 - (x_{2} - \mu_{2})^2 \nonumber 
\end{eqnarray} \pause

Call $\boldsymbol{\mu} = (\mu_{1}, \mu_{2} ) $\\ \pause

The Hessian at the critical value is \pause
\begin{eqnarray}
\boldsymbol{H}(f)(\boldsymbol{\mu} ) & = & \begin{pmatrix} 
\frac{\partial^{2} U_{i}} {\partial x_{1} \partial x_{1}} (\boldsymbol{\mu} ) & \frac{\partial^{2} U_{i}} {\partial x_{1} \partial x_{2}} (\boldsymbol{\mu} )  \\
\frac{\partial^{2} U_{i}} {\partial x_{2} \partial x_{1}} (\boldsymbol{\mu} ) & \frac{\partial^{2} U_{i}} {\partial x_{2} \partial x_{2}} (\boldsymbol{\mu} )\\
\end{pmatrix} \nonumber \\
& = & \begin{pmatrix} 
-2 & 0  \\
0 & -2  \\
\end{pmatrix} 
\nonumber 
\end{eqnarray} \pause
So, $-2 * -2 - 0  = 4 >0$ and $-2 <0$ $\leadsto$ \alert{negative definite}, \alert{maximum} \\
$\boldsymbol{\mu} = (\mu_{1}, \mu_{2})$ are legislator $i$'s two dimensional ideal point.  

\end{frame}








\end{document}
